% from ega4-2-fr.tex
\oldpage[IV-2]{193}

Si $X$, $Y$ sont deux préschémas localement noethériens, nous dirons qu'un
morphisme $f:X\to Y$ est un \emph{$\boldsymbol P$-morphisme} si :
\begin{enumerate}
\item[1\textsuperscript{o}] $f$ est \emph{plat};
\item[2\textsuperscript{o}] pour tout $y\in Y$, la propriété
$\boldsymbol P(f^{-1}(y),k(y))$ est vraie.
\end{enumerate}

\begin{lemma}[7.3.2]
\label{IV.2.s7.3.2-fr}
Soient $X$, $Y$ deux préschémas localement noethériens,
$f:X\to Y$ un morphisme. Les propriétés suivantes sont équivalentes :
\begin{enumerate}
\item[\emph{a)}] $f$ est un $\boldsymbol P$-morphisme.
\item[\emph{b)}] Pour tout $x\in X$, si l'on pose $y=f(x)$, alors $\mathscr O_x$ est
un $\mathscr O_y$-module plat et la propriété
$\boldsymbol Q(\mathscr O_x\otimes_{\mathscr O_y}k(y),k(y))$ est vraie.
\item[\emph{c)}] Pour tout $x\in X$, si l'on pose $y=f(x)$, le morphisme
$\operatorname{Spec}(\mathscr O_x)\to\operatorname{Spec}(\mathscr O_y)$
correspondant à l'homomorphisme $\mathscr O_y\to\mathscr O_x$ est un
$\boldsymbol P$-morphisme.
\item[\emph{c$'$)}] Pour tout point fermé $x\in X$, le morphisme
$\operatorname{Spec}(\mathscr O_x)\to\operatorname{Spec}(\mathscr O_y)$
est un $\boldsymbol P$-morphisme.
\end{enumerate}
\end{lemma}

\begin{proof}
L'équivalence de a) et b) résulte en effet des définitions; il en est de
même de l'équivalence de b) et c), compte tenu de (I, 2.4.2); enfin
l'équivalence de b) et c$'$) résulte de ce que pour tout $x\in X$,
$\overline{\{x\}}$ contient un point fermé (5.1.11).
\end{proof}

\begin{corollary}[7.3.3]
\label{IV.2.s7.3.3-fr}
Soient $A$, $B$ deux anneaux noethériens, $\varphi:A\to B$ un homomorphisme
tel que le morphisme correspondant
$\operatorname{Spec}(B)\to\operatorname{Spec}(A)$ soit un
$\boldsymbol P$-morphisme. Alors, pour toute partie multiplicative $S$ de
$A$, le morphisme
$\operatorname{Spec}(S^{-1}B)\to\operatorname{Spec}(S^{-1}A)$ est un
$\boldsymbol P$-morphisme.
\end{corollary}

\begin{proof}
Cela résulte aussitôt de (7.3.2) et de (I, 1.6.2).
\end{proof}

\begin{env}[7.3.4]
\label{IV.2.s7.3.4-fr}
Nous supposerons toujours dans ce qui suit que la propriété
$\boldsymbol Q$ est telle que les trois conditions suivantes soient
vérifiées :
\begin{enumerate}
\item[$(\mathrm P_{\mathrm I})$] \emph{(transitivité).} --- Si
$f:X\to Y$ est un morphisme régulier (6.8.1) et $g:Y\to Z$ un
$\boldsymbol P$-morphisme, alors $g\circ f$ est un
$\boldsymbol P$-morphisme.
\item[$(\mathrm P_{\mathrm{II}})$] \emph{(descente).} --- Si
$f:X\to Y$ et $g:Y\to Z$ sont des morphismes de préschémas localement
noethériens tels que $f$ soit fidèlement plat et que $g\circ f$ soit un
$\boldsymbol P$-morphisme, alors $g$ est un $\boldsymbol P$-morphisme.
\item[$(\mathrm P_{\mathrm{III}})$] --- Pour tout corps $k$, la propriété
$\boldsymbol P(\operatorname{Spec}(k),k)$ est vraie.
\end{enumerate}
\end{env}

\begin{remarks}[7.3.5]
\label{IV.2.s7.3.5-fr}
\begin{enumerate}
\item[\emph{i)}] Les conditions $(\mathrm P_{\mathrm I})$ et
$(\mathrm P_{\mathrm{III}})$ entraînent que \emph{tout morphisme régulier
est un $\boldsymbol P$-morphisme}.
\item[\emph{ii)}] Notons que les hypothèses de $(\mathrm P_{\mathrm I})$
(resp. $(\mathrm P_{\mathrm{II}})$) entraînent que $h=g\circ f$ (resp.
$g$) est plat (2.2.13); les hypothèses de $(\mathrm P_{\mathrm I})$ ou de
$(\mathrm P_{\mathrm{II}})$ entraînent d'autre part que pour tout
$z\in Z$, $f_z:h^{-1}(z)\to g^{-1}(z)$ est plat (2.1.4); comme, pour tout
$y\in g^{-1}(z)$, $f_z^{-1}(y)$ est isomorphe à $f^{-1}(y)$ (I, 3.6.4),
cela montre qu'il suffit, pour vérifier les conditions
$(\mathrm P_{\mathrm I})$ et $(\mathrm P_{\mathrm{II}})$, de le faire
seulement lorsque $Z$ est le \emph{spectre d'un corps}.
\item[\emph{iii)}] Dans certains cas, la propriété $\boldsymbol Q$ sera telle que la
condition suivante soit vérifiée :
\begin{enumerate}
\item[$(\mathrm P'_{\mathrm I})$] --- Si $f:X\to Y$ et $g:Y\to Z$ sont
deux $\boldsymbol P$-morphismes, alors $g\circ f$ est un
$\boldsymbol P$-morphisme.
\end{enumerate}
\end{enumerate}
\end{remarks}

